%! @@TITLE@@ — Unit @@UNITINT@@, Lesson @@LESSONID@@ — Lesson Plan
% GRADUAL-RELEASE plan (warm-up / guided notes and practice / group activity / debrief /
% close and assign). Teacher-facing — use the formal AP vocabulary freely.
% NOTE: this course has no `fixedskillbox`. Every lesson-plan box is `skillbox`, which is
% breakable; put \boxguard (raised, e.g. \boxguard[30], ahead of a tabularx that cannot
% split) before each one instead. There are no tier boxes and no exit ticket.
\documentclass[10pt]{article}
\usepackage{@@PREFIX@@-article}
\usepackage{@@PREFIX@@-boxes}
\usepackage{graphicx}   % -article does NOT load graphicx; needed for thumbnails/screenshots
\graphicspath{{images/}}
\newcommand{\TallMath}[1]{$\displaystyle #1\rule[-1.4em]{0pt}{3.2em}$}
@@COURSEMACROS@@\newcommand{\UnitNumberName}{Unit @@UNITINT@@: @@UNITTITLE@@ \quad}
\newcommand{\LessonNumberName}{Lesson @@LESSONID@@: @@TITLE@@}

\begin{document}
\begin{center}
    {\Large\bfseries \CourseName \\
    {\small\bfseries \UnitNumberName \LessonNumberName}}
\end{center}

\begin{tcolorbox}[colback=sky, colframe=navy, boxrule=0.9pt, arc=2mm,
    left=2mm, right=2mm, top=1.5mm, bottom=1.5mm]
    \textbf{Primary Objective:} TODO --- what students can do/interpret/justify by the end.
    \\[2pt]
    \textbf{CED alignment:} Topic TODO \quad$\cdot$\quad Big Idea: TODO (VAR / UNC / DAT / PROB)
    \quad$\cdot$\quad Skill TODO \quad$\cdot$\quad LO TODO \quad$\cdot$\quad EK TODO
    \\[2pt]
    \textbf{Lesson model:} \emph{Gradual release --- I do, we do, you do.} The guided notes
    deliver the vocabulary and the core idea directly; the guided practice is worked together with
    students holding the pen; the independent practice and the group activity are where students
    carry it alone and then with peers. The debrief closes the loop as a whole class.
\end{tcolorbox}

\renewcommand{\arraystretch}{1.4}
\begin{skillbox}[Learning Targets \& Key Understandings]{goldbox}
% Fill in the \item text ONLY. Two tabularx cells, one itemize each — do not re-nest or
% re-balance the environments (that is the classic source of a stray \end{itemize}).
\begin{tabularx}{\linewidth}{@{}X|X@{}}
\textbf{Learning targets (student language)} & \textbf{Key understandings (the ``why'')} \\[2pt]
\hline
\vspace{-6pt}
\begin{itemize}[leftmargin=1.1em, nosep, after=\vspace{2pt}]
  \item TODO: ``I can \dots'' target, traceable to the CED Learning Objective.
\end{itemize}
&
\vspace{-6pt}
\begin{itemize}[leftmargin=1.1em, nosep, after=\vspace{2pt}]
  \item TODO: the why (paraphrase the Essential Knowledge) --- and state the lesson's
        \textbf{target misconception} explicitly.
\end{itemize}
\end{tabularx}
\end{skillbox}

\begin{skillbox}[{Vocabulary, Concepts, \& Theorems --- taught directly in the guided notes}]{greenbox}
    % TODO: key terms + definitions for this lesson. Two \begin{minipage}[t]{0.48\linewidth}
    % columns of \itemize, separated by \hfill.
\end{skillbox}

% A tabularx never splits, so guard generously ahead of it.
\boxguard[30]
\begin{skillbox}[Lesson at a Glance --- \MeetingLength]{sky}
\renewcommand{\arraystretch}{1.25}
\begin{tabularx}{\linewidth}{@{}l c X X@{}}
\textbf{Phase} & \textbf{Min} & \textbf{Students} & \textbf{Teacher} \\
\hline
Warm-Up & 5 & TODO & Circulate; debrief the seed item aloud and leave it on the board \\
Guided Notes \& Practice & 22 & Fill the vocabulary box and notes sections; then TODO together; then TODO alone & \textbf{I do} the notes $\to$ \textbf{we do} the Guided Practice box $\to$ \textbf{you do} Independent Practice while you circulate \\
Group Activity & 16 & TODO activity title, scenarios 1--2, groups of 3--4 & Circulate; press for evidence from the data; choose group work to display in the debrief \\
Debrief & 8 & Groups report out; fill the Debrief box together & Take the crux public: correct the near-miss answer gently and in front of everyone \\
Close \& Assign & 4 & Start the homework; hear the preview & Announce the homework (or the DeltaMath substitute) and preview the next lesson \\
\end{tabularx}
\end{skillbox}

\boxguard[20]
\begin{skillbox}[Warm-Up --- Activate Prior Knowledge (5 min)]{sky}
    % TODO: the ~3 warm-up items and the SEED each plants (which part of the notes it sets up),
    % plus what to debrief aloud and what to leave on the board for the notes to point back at.
@@SPIRALWARMUP@@
\end{skillbox}

\boxguard
\begin{skillbox}[Guided Notes \& Practice --- I do, we do, you do (22 min)]{sky}
\begin{multicols}{2}
\textbf{I do (about 12 min).} Present the notes sections in order; students fill the blanks as
each idea is named, and the vocabulary box is filled \emph{as you go}, never in advance. TODO:
which section is the must-land one and which can be compressed; where the pre-drawn display goes
on the board; how the target misconception is surfaced with evidence rather than assertion.

\columnbreak
\textbf{We do (about 6 min).} The Guided Practice box, \emph{TODO Title}. Students hold the pen;
you hold the questions: TODO the three questions you ask. \textbf{Part TODO is the point of the
box} --- it is the one that tests the misconception on new ground.

\textbf{You do (about 4 min).} Independent Practice, three items, silent and alone. Circulate.
\textbf{Item~TODO is the formative check} --- glance at it on every desk. If more than a third
TODO, reteach before releasing the activity and cut activity scenario~1 to items (a)--(b) to buy
the time.
\end{multicols}
\end{skillbox}

\boxguard
\begin{skillbox}[Group Activity --- \emph{TODO Title} (16 min)]{sky}
\textbf{One activity for the whole class --- there is no tiered version.} Groups of 3--4.
\textbf{Launch (1 min).} TODO: the launch script --- the data context and the ground rule
(\emph{nobody writes an answer until somebody can point at the data and say how they know}).

\begin{multicols}{2}
\textbf{What students do.} TODO: what scenario~1 applies (the clean case, so every group has
something correct before the crux), what scenario~2 varies, and which item is the
\textbf{crux question} (the target misconception).

\columnbreak
\textbf{What the teacher does.} Circulate and press for evidence from the data rather than
supplying answers:
\begin{itemize}[leftmargin=1.1em, nosep]
  \item TODO: item ref: ``TODO prompt.''
\end{itemize}
While circulating, pick the work to display in the debrief --- including one \emph{almost}-right
answer to the crux.
\end{multicols}
\end{skillbox}

\boxguard
\begin{skillbox}[Debrief --- whole class (8 min)]{sky}
Groups report out, then fill the Debrief box at the end of the activity together.
\begin{multicols}{2}
\begin{enumerate}[leftmargin=1.3em, nosep, itemsep=3pt]
  \item \textbf{Open on the \emph{almost}-right crux answer} you collected. Correcting it in
        public, gently, is what makes the crux stick. TODO: how you stage it.
  \item TODO: the next report-out to take, and what to write beside it.
\end{enumerate}
\columnbreak
\textbf{The Debrief box.} Three questions, filled together as the report-out lands. TODO: what
each one is for --- item~1 is the sentence you want on every desk.

\textbf{If the debrief runs short}, cut TODO --- Debrief item~1 is the only one that cannot
be cut.
\end{multicols}
\end{skillbox}

\boxguard
\begin{skillbox}[AP Practice --- assigned, not scored]{goldbox}
TODO: four AP-style multiple-choice items (five options each, in context) plus one multi-part
free-response set, and what each targets. Name the spiral item.

\textbf{Not scored --- the cover's score column reads \textbf{NA} for this page.} Go over it at
the start of the next lesson and hold the AP standard out loud: \emph{in context or it does not
count}. TODO: the part worth reading closely, and what does not earn the point.
\end{skillbox}

\boxguard
\begin{skillbox}[Homework --- scored, due next class]{goldbox}
TODO: the assignment, in fresh contexts, and what each item targets. \textbf{Start it in the last
four minutes of the period} so students hit item~1 while you are still in the room.
\textbf{Scoring:} TODO points --- TODO breakdown. \textbf{Item~TODO is the one to read closely}:
TODO why it predicts performance on the next lesson.

\textbf{DeltaMath override:} when DeltaMath carries this topic, assign it in place of this page,
strike through the score cell on the cover, and leave this page in the packet as extra practice.
\end{skillbox}

\boxguard
\begin{skillbox}[Watch For (while circulating)]{redbox}
\begin{itemize}[leftmargin=1.2em, nosep]
  \item TODO: misconception, keyed to an item number.
\end{itemize}
Cold-call prompts: TODO.
\end{skillbox}

\boxguard
\begin{skillbox}[Close \& Assign (4 min)]{goldbox}
Hand the homework launch first --- TODO, scored, due next class (or the DeltaMath assignment in
its place) --- and point out that the AP Practice page is theirs to work and bring to review.
Then close on what changed rather than on the assignment: TODO one sentence.
\textbf{Preview:} TODO next lesson.
\end{skillbox}

% ── Teacher notes — one per component, in packet order ──
% TEACHERNOTES: teacher prose lives HERE, never in a _key. A note in a key is the one block
% with no counterpart in the blank, so it makes the key run longer and costs the student
% packet a blank page.
\begin{teachernote}[Warm-Up]
TODO: pacing + what each item seeds; what to cut if it runs long; what to leave on the board.
\end{teachernote}

\begin{teachernote}[Guided Notes \& Practice]
\textbf{Pace it 12 / 6 / 4.} TODO: which sections are must-land and which can be compressed.
The vocabulary box is filled \textbf{as each term is used}, never front-loaded.
\textbf{Guided Practice is where the pen changes hands} --- take answers, write what students
say, and push back rather than working it on the board while they watch.
\textbf{Independent Practice is the formative check.} TODO: the item to read over shoulders, the
three piles to sort responses into, and the reteach trigger.
\end{teachernote}

\begin{teachernote}[Group Activity]
\textbf{16 minutes, roughly TODO on scenario~1 and TODO on scenario~2.} \textbf{One activity for
the whole class:} do not split the room by level --- the differentiation is in how you circulate.
TODO: the early-finisher move.
\textbf{Debrief (8 min).} TODO: the must-land moments; what to cut if time is short.
\end{teachernote}

\begin{teachernote}[AP Practice]
\textbf{This page is not scored} --- the graded practice is the homework that follows it. TODO:
what each item targets; the part worth reading closely; when to go over it.
\end{teachernote}

\begin{teachernote}[Homework]
\textbf{Scored, TODO points, due next class.} Launch it in the last four minutes.
TODO: the item that predicts next-lesson performance and why; what does not earn the point.
\textbf{When DeltaMath carries this topic, assign it instead}, strike the score cell on the
cover, and tell students this page is now optional practice. Do not assign both.
\end{teachernote}

\end{document}
